% ============================================================================
% CHAPITRE 7 -- ASSISTANT CONVERSATIONNEL RAG
% Sections « sprint » : introduction, objectif, backlog
% A inserer avant la section « Contexte et besoin »
% ============================================================================

\section{Introduction}

Ce chapitre présente le second volet du sprint~4, consacré à l'assistant
conversationnel réglementaire. Le chapitre précédent a traité le premier volet,
le tableau de bord décisionnel destiné à l'administrateur~; celui-ci s'adresse
aux usagers des espaces forestiers et répond au manque d'accès à une information
réglementaire adaptée, relevé dès la critique de l'existant.

Après avoir exposé l'objectif et le backlog de ce volet, le chapitre présente les
fondements technologiques retenus, le principe général d'un pipeline de génération
augmentée par récupération (RAG), puis sa réalisation concrète dans le microservice
\texttt{chatbot\_ms}.

\section{Objectif du sprint}

Ce volet vise à mettre à disposition des usagers de la forêt un assistant capable
de répondre à leurs questions réglementaires à partir des textes officiels
tunisiens, plutôt que des connaissances internes d'un modèle de langage.

Les principaux objectifs sont les suivants~:

\begin{itemize}
  \item \textbf{Constituer une base documentaire réglementaire} organisée par
        spécialité (chasse, camping, apiculture) à partir du Code forestier et
        du Journal Officiel de la République Tunisienne, complétée d'une aide à
        l'usage de l'application.
  \item \textbf{Mettre en place un pipeline RAG} combinant recherche sémantique
        et recherche lexicale, afin de retrouver les passages pertinents avant
        toute génération de réponse.
  \item \textbf{Garantir l'étanchéité entre spécialités}~: un usager ne reçoit
        jamais un extrait relevant d'une activité autre que la sienne.
  \item \textbf{Encadrer la génération} de manière que chaque affirmation soit
        fondée sur un extrait récupéré et accompagnée de sa source, réduisant
        ainsi le risque de réponse non fondée.
  \item \textbf{Déployer un microservice indépendant} (\texttt{chatbot\_ms}),
        sans base de données relationnelle, dont l'index est stocké sous forme
        de fichiers montés en volume Docker.
\end{itemize}

À l'issue de ce sprint, un usager dont le dossier a été approuvé peut interroger
l'assistant depuis l'application qui lui est dédiée et obtenir une réponse
contextualisée, limitée à sa spécialité et accompagnée des articles sur lesquels
elle s'appuie.

\section{Backlog du sprint}

Le tableau~\ref{tab:backlog-sprint4-chatbot} présente le backlog de ce volet. Il
correspond à la User Story~\#30, unique élément de l'EPIC~9 du product backlog,
décomposée en tâches techniques.

\begin{table}[h]
\centering
\caption{Backlog du sprint 4 --- volet assistant conversationnel}
\label{tab:backlog-sprint4-chatbot}
\begin{tabular}{|c|p{3cm}|p{7cm}|c|c|}
\hline
\textbf{ID} & \textbf{Titre} & \textbf{Tâches techniques} & \textbf{Priorité} & \textbf{SP} \
\hline
30 & Chatbot réglementaire scopé par spécialité &
\begin{minipage}[t]{7cm}
\vspace{2pt}
\begin{itemize}[leftmargin=*,nosep]
  \item Constituer le corpus par domaine (\texttt{chasse}, \texttt{camper},
        \texttt{apiculture}, \texttt{app}) et documenter chaque fichier par un
        \texttt{\_manifest.json} (source, rang normatif, validité)
  \item Implémenter le chunking juridique par article
        (\texttt{rag/chunker.py})~: décomposition des tableaux ligne à ligne et
        redécoupage des articles trop longs
  \item Construire l'index hors-ligne
        (\texttt{scripts/build\_index.py})~: embeddings \texttt{BAAI/bge-m3} et
        index lexical BM25
  \item Développer la recherche hybride~: recherche dense par similarité
        cosinus, recherche BM25 avec tokenisation et expansion de requête,
        fusion des classements par RRF
  \item Filtrer les résultats par domaine selon la spécialité de l'usager
        extraite du jeton JWT
  \item Encadrer la génération via le LLM Groq~: prompt système imposant la
        fidélité aux extraits, la citation de l'article et la hiérarchie des
        normes
  \item Exposer \texttt{POST /api/chat/} et \texttt{/api/chat/debug}, charger
        modèle et index au démarrage du service, et mettre en place un cache
        de réponses sur disque
  \item Intégrer l'écran de conversation dans \texttt{public\_user\_app}
\end{itemize}
\vspace{2pt}
\end{minipage}
& \textcolor{orange}{\textbf{Haute}} & 8 \
\hline
\end{tabular}
\end{table}

\textbf{Note de périmètre}~: l'EPIC~9 ne comporte qu'une seule User Story, celle-ci
couvrant à elle seule l'assistant conversationnel. Le périmètre réel du sprint est
donc porté par sa décomposition en tâches techniques, selon la même démarche que
celle retenue pour les sprints précédents.
